\begin{pastexam}[label={px:gw3-102-ds-04}]{關務三等資訊處理 102 年 第四題}
（申論；30 分）

考慮排序（sort）的問題：（每小題 10 分，共 30 分）
\begin{itemize}[leftmargin=2.2em]
\item[㈠] 如果要排序的資料很少，例如只有十幾筆資料，那麼你將採用快速排序法（quick sort）？合併排序法（merge sort）？
還是氣泡排序法（bubble sort）？為什麼？
\item[㈡] 如果要排序的資料很多，例如多到超過主記憶體容量許多，那麼你將採用快速排序法？合併排序法？還是氣泡排序法？
為什麼？
\item[㈢] 快速排序法、合併排序法以及氣泡排序法這三個排序法當中，那一（些）排序法是穩定的（stable）？或者都不穩定？
\end{itemize}
\end{pastexam}

\begin{pastsoln}
答案：㈠ 氣泡排序法；㈡ 合併排序法；㈢ 合併排序法和氣泡排序法穩定，快速排序法不穩定。 ㈡ 是外部排序（資結書
第~15.3~節），整題以第~14~章的內容為主。

\textbf{㈠ 氣泡排序法}。 漸近複雜度只描述 $n$ 很大時的趨勢，$n$ 只有十幾時，常數和額外的固定成本才是重點：
\begin{itemize}
\item 快速排序、合併排序的 $O(n \log n)$ 要靠遞迴達成，每層呼叫有參數傳遞、堆疊的成本；合併排序還要配置
暫存陣列、來回複製。 $n$ 小時 $n \log n$ 和 $n^2$ 差不了多少（$n = 15$ 時 $n^2/2 \approx 112$，
$n \lg n \approx 59$），這些額外成本反而佔大部分。
\item 氣泡排序只有兩層迴圈和相鄰交換，程式短、不用額外空間、不容易寫錯；加上「一趟沒有交換就停」的檢查，資料
大致排好時接近 $O(n)$。
\end{itemize}
所以資料很少時選氣泡排序。 實務上的函式庫正是這麼做的：快速排序或合併排序遞迴到子陣列只剩十幾個時，改用
簡單的平方時間排序（通常是插入排序，和氣泡排序同類但比較次數更少）收尾。

\textbf{㈡ 合併排序法}。 資料超過主記憶體很多時，必須放在磁碟上排序（外部排序），這時成本主要是磁碟 I/O 的次數，
而磁碟適合\textbf{循序}讀寫，不適合隨機存取：
\begin{itemize}
\item 合併排序天生是循序的：先每次讀入記憶體放得下的一段，在記憶體內排好後寫回磁碟，成為一個個已排序的
「串」（run）；再把好幾個串同時從頭到尾循序讀入、合併、循序寫出。 $r$ 個初始串用 $k$ 路合併時只要
$\lceil \log_k r \rceil$ 趟，每趟把全部資料讀寫一次（資結書第~15.3~節）。
\item 快速排序的分割要從兩端往中間掃、交換相隔很遠的元素，放在磁碟上就是大量隨機存取，而且遞迴分割的趟數
是 $O(\log n)$ 到 $O(n)$，不適合。
\item 氣泡排序要 $O(n^2)$ 次比較、交換，資料量大時完全不可行。
\end{itemize}

\textbf{㈢} \textbf{合併排序和氣泡排序是穩定的，快速排序不穩定}（資結書第~14.7~節）。
\begin{itemize}
\item 氣泡排序只在相鄰兩筆「嚴格大於」時才交換，相等的不交換，相等資料的相對順序不變。
\item 合併排序合併時，左右兩段的首筆相等就先取左段的（比較寫成 $\le$），左段的資料原本就在前面，所以穩定。
（若寫成 $L[i] < R[j]$ 才取左段，相等時會先取右段，就不穩定了。）
\item 快速排序分割時會把元素和遠處的元素交換，跳過中間相等的資料。 例：$\langle 2_a, 2_b, 1\rangle$ 以最後一個 $1$
為基準做 Lomuto 分割，沒有元素 $\le 1$，最後把基準換到最前面，$1$ 和 $2_a$ 交換，得到
$\langle 1, 2_b, 2_a\rangle$，兩個 $2$ 的順序反了。
\end{itemize}
\end{pastsoln}
